\section{Metodolog\'ia experimental y validaci\'on}
\label{sec:metodologia}

\subsection{Estrategia general de validaci\'on}

La validaci\'on del sistema se organiz\'o en cuatro niveles, de menor a mayor integraci\'on:

\begin{enumerate}[itemsep=0.15em]
    \item \textbf{Verificaci\'on de modelos} contra medici\'on de laboratorio (modelo del
    SM-24, modelo del amplificador operacional real, ganancias de continua de cada etapa).
    \item \textbf{Pruebas unitarias y de regresi\'on automatizadas} sobre el firmware y el
    software de procesamiento.
    \item \textbf{Caracterizaci\'on cuantitativa en banco} con instrumental de referencia
    (osciloscopio de cuatro canales) sobre m\'ultiples placas y configuraciones.
    \item \textbf{Pruebas de aceptaci\'on de sistema completo} (auditor\'ia pre-campo) y
    campa\~nas de medici\'on en campo.
\end{enumerate}

\subsection{Modelado y caracterizaci\'on del transductor}

El ge\'ofono SM-24 se model\'o en Simulink a partir de las ecuaciones de la
Secci\'on~\ref{sec:geo-tres-formas}, y el
modelo se contrast\'o con mediciones sobre el transductor real. Se desarrollaron
herramientas propias de an\'alisis:

\begin{itemize}[itemsep=0.15em]
    \item \texttt{geophone\_scope.m} y \texttt{geophone\_scope\_simple.m}: adquisici\'on y
    visualizaci\'on en tiempo real sobre puerto serie.
    \item \texttt{analisis\_espectral.m}: estimaci\'on de densidad espectral de potencia
    por el m\'etodo de Welch y c\'alculo de relaci\'on se\~nal--ruido por entidad de medici\'on.
    \item \texttt{ComparacionFiltros.m}: comparaci\'on cuantitativa de topolog\'ias de
    filtro, que fundament\'o la elecci\'on de MFB.
    \item Protocolo de barridos de la cadena anal\'ogica con \emph{script} de an\'alisis
    dedicado, que permiti\'o contrastar la respuesta medida contra el modelo del circuito
    con valores de componentes reales.
\end{itemize}

\begin{sloppypar}
Para hacer trazable el estado MASW mostrado en resultados, se fija el siguiente
\emph{snapshot}: \texttt{data} en el commit
{\ttfamily\footnotesize 76c4254b19cf\allowbreak c1f85edb9fb9\allowbreak
201b40c7f250\allowbreak b413},
software de procesamiento en
{\ttfamily\footnotesize 30b60517d337\allowbreak 7d175ecc7070\allowbreak
278c25cbca35\allowbreak b7ec}
y archivo
\nolinkurl{data/processed/Canchita/field_review_masw_state.json}, con SHA-256
{\ttfamily\footnotesize 8f8bf4d59a09\allowbreak 2b8f321b2829\allowbreak
af07fbfd3e03\allowbreak 374304c57310\allowbreak 3c39e1a43df3\allowbreak 0006}.
El estado guarda: imagen $c=\SIrange{50}{170}{\meter\per\second}$ con paso
\SI{0.1}{\meter\per\second} y $f=\SIrange{0.1}{35}{\hertz}$; grupo~1 con peso~1 y
grupo~2 con peso~0; 113 puntos etiquetados como modo~0; e inversi\'on de siete capas
m\'as semiespacio, 1000 iteraciones, $b_s=8$, $b_h=12$, $\nu=0{,}35$ y
$\rho=\SI{1850}{\kilogram\per\meter\cubed}$. El desajuste del \emph{port} se define como
\end{sloppypar}

\begin{equation}
E=\frac{100}{n}\sum_{i=1}^{n}\frac{|c_i^{\mathrm{obs}}-c_i^{\mathrm{teo}}|}
{c_i^{\mathrm{obs}}}\quad[\%].
\end{equation}

Este \emph{snapshot} permite repetir el estado num\'erico guardado, pero no reconstruye
por s\'i solo todas las decisiones manuales que llevaron al pol\'igono de selecci\'on. Esa
limitaci\'on es una raz\'on adicional para tratar la inversi\'on como exploratoria. La SNR
de la Fig.~\ref{fig:gather-main}, por su parte, se calcul\'o sobre la primera captura
v\'alida a \SI{1020}{\hertz} de cada posici\'on, no sobre una traza apilada: cociente entre
la desviaci\'on est\'andar posterior y previa al impacto, expresado como
$20\log_{10}(\sigma_{\mathrm{s}}/\sigma_{\mathrm{r}})$, donde los sub\'indices
$\mathrm{s}$ y $\mathrm{r}$ identifican las ventanas de se\~nal y ruido.

El modelo ideal no reprodujo el comportamiento medido del compensador en el extremo
superior de la banda. El primer modelo no ideal usado en los c\'alculos combin\'o
$A_0=\SI{90}{\decibel}$ con \SI{8}{\mega\hertz}, pero confundi\'o frecuencia de ganancia
unitaria con polo dominante. Esa parametrizaci\'on queda invalidada y debe recalcularse;
la funci\'on monopolo ya fue corregida en los scripts, pero la s\'intesis y la medici\'on
deben repetirse. El resultado anterior no se usa para sostener una banda compensada ya
demostrada.

\begin{figure}[ht]
\centering
\includegraphics[width=\textwidth]{barrido_crudo_cadena.png}
\caption{Ejemplo de barrido experimental conservado como evidencia de laboratorio:
se\~nales crudas de excitaci\'on y salida de la cadena antes de la estimaci\'on arm\'onica.
Los puntos medidos siguen siendo v\'alidos como registro; cualquier superposici\'on con
el modelo no ideal anterior debe recalcularse con el polo dominante corregido.}
\label{fig:barrido-crudo-metodo}
\end{figure}

\subsection{Validaci\'on experimental del DA-AFE}
\label{sec:met-daafe}

Este es el experimento m\'as riguroso realizado hasta la fecha y fue documentado en la
comunicaci\'on sometida a URUCON~2026; a\'un no se cuenta con dictamen externo.

\subsubsection{Dise\~no experimental}

Se evaluaron \textbf{ocho configuraciones placa--ganancia}: las cinco ganancias del PGA
sobre la Placa~A, m\'as las ganancias $\times16$, $\times32$ y $\times50$ sobre la
Placa~B. Cada configuraci\'on se evalu\'o mediante \textbf{diez reinicios apareados}
($N=10$).

Las decisiones metodol\'ogicas relevantes fueron:

\begin{itemize}[itemsep=0.15em]
    \item \textbf{Referencia com\'un}: se aplic\'o una \'unica l\'inea base de c\'odigos fijos,
    elegida emp\'iricamente sobre la Placa~B a $\times50$, \emph{sin} adaptarla a cada
    configuraci\'on. Esto es deliberadamente desfavorable a la l\'inea base: las reducciones
    reportadas acotan superiormente el beneficio respecto de un ajuste fijo por placa, en
    lugar de replicarlo.
    \item \textbf{Configuraci\'on \'unica de controlador}: los mismos par\'ametros de la
    Tabla~\ref{tab:cal_params} para todas las placas y ganancias, sin reajuste.
    \item \textbf{Medici\'on apareada}: dentro de cada reinicio se registr\'o siempre primero
    la condici\'on de referencia fija y luego la calibrada.
    \item \textbf{Escritura en EEPROM deshabilitada} tras la calibraci\'on, de modo que
    cada reinicio partiera exactamente de la misma l\'inea base.
    \item \textbf{Exclusiones declaradas}: la Placa~B a $\times1$ y $\times4$ se
    excluyeron porque sus residuos de l\'inea base ya ca\'ian dentro de las bandas muertas,
    impidiendo actualizaciones correctivas. El comportamiento esperado se confirm\'o con
    un diagn\'ostico separado de banda muerta m\'as estrecha.
\end{itemize}

\subsubsection{Procedimiento de medici\'on}

Un osciloscopio Tektronix TDS~2004B de cuatro canales registr\'o la tensi\'on media de
salida de las etapas PGA, BP, SUM y LP sobre una ventana de \SI{1}{\second} posterior al
establecimiento. Las masas de las puntas se conectaron a $V_{\mathrm{ref}}=V_{DDA}/2$,
de modo que cada canal midiera $V_{\mathrm{etapa}}-V_{\mathrm{ref}}$. La salida LP se
adopt\'o como m\'etrica principal por alimentar directamente al ADC.

\subsubsection{An\'alisis estad\'istico}

Se aplic\'o una \textbf{prueba de signos de una cola}~\cite{Conover1999} sobre las
mediciones apareadas de la etapa LP. La elecci\'on de una prueba no param\'etrica responde
al tama\~no de muestra y a no asumir normalidad de los residuos.

\subsection{Auditor\'ia pre-campo del sistema completo}
\label{sec:auditoria}

Antes de las campa\~nas de campo se ejecut\'o una auditor\'ia sistem\'atica sobre hardware
real, con una matriz de pruebas de aceptaci\'on identificadas E1--E18 que cubri\'o: capturas
en RAM y en SD de distinta duraci\'on, decimaci\'on y filtrado por radio, comandos de
captura a SD, verificaci\'on estricta por USB, interfaz en navegador, reconexi\'on de
WebSocket, detenci\'on en vivo, transiciones de rango y decimaci\'on, ciclo de apagado y
recuperaci\'on, y prueba de resistencia prolongada.

\subsubsection{Defectos cr\'iticos encontrados y corregidos}

La auditor\'ia detect\'o defectos que habr\'ian invalidado una campa\~na de campo. Se destacan
tres por su valor metodol\'ogico:

\begin{itemize}[itemsep=0.25em]
    \item \textbf{F1 --- Capturas largas destruidas por el \emph{watchdog}.} Un
    \emph{watchdog} de estancamiento de sincronizaci\'on re-disparaba el comando de arranque
    y destru\'ia las sesiones largas en SD. El s\'intoma observado era desconcertante: una
    captura de 10 minutos aceptada con \emph{cero muestras}. Se corrigi\'o con un margen
    proporcional y un aborto limpio.
    \item \textbf{F9 --- El arranque romp\'ia el nodo.} Defecto de carrera:
    el maestro esperaba la confirmaci\'on de armado con un tiempo l\'imite de
    \SI{120}{\milli\second}, pero la confirmaci\'on real de ``PSoC armado'' tarda
    \SI{130}{\milli\second}. El reintento reemit\'ia la secuencia completa de preparaci\'on,
    forzando un segundo comando de configuraci\'on; pero el PSoC, ya armado, entra en
    silencio total por dise\~no. El comando nunca se confirmaba, y el esclavo interpretaba
    el vencimiento como fallo y liberaba el almac\'en, mientras el PSoC segu\'ia f\'isicamente
    armado esperando un flanco de sincronismo que nadie volver\'ia a emitir. El nodo quedaba
    inutilizable hasta un reinicio f\'isico. Es un ejemplo did\'actico de c\'omo una carrera
    de \SI{10}{\milli\second} entre dos temporizadores produce una falla determinista y
    total. Se corrigi\'o en dos capas: el reintento dej\'o de reemitir la preparaci\'on, y el
    comando de preparaci\'on se hizo idempotente.
    \item \textbf{F5 --- Cliente colgado en conexi\'on semiabierta.} La herramienta de
    descarga no ten\'ia \emph{watchdog} de silencio y quedaba bloqueada sobre un
    \emph{socket} zombi.
\end{itemize}

\subsubsection{Criterio de aceptaci\'on}

La prueba de aceptaci\'on definitiva consisti\'o en una captura continua de 10 minutos con
verificaci\'on de integridad de extremo a extremo: conteo de lotes en SD, conteo de
muestras recibidas por radio, tama\~no en bytes, n\'umero de conexiones y desconexiones, y
recomputaci\'on de la suma SHA-256. La adopci\'on de un criterio de integridad
criptogr\'afica ---en lugar de una inspecci\'on visual de las se\~nales--- fue deliberada:
permite distinguir una transferencia completa de una truncada, distinci\'on que result\'o
decisiva, ya que un intento previo hab\'ia capturado correctamente en SD pero entregado un
archivo truncado por el lado del cliente.

\subsection{El sitio de ensayo: contexto geol\'ogico y verdad de terreno disponible}
\label{sec:sitio-sev}

Un elemento poco habitual en un trabajo de instrumentaci\'on, y que conviene explotar, es
que el sitio principal de medici\'on cuenta con una \textbf{caracterizaci\'on geof\'isica
independiente, previa y co-localizada}, realizada por terceros con un m\'etodo distinto.

En julio de 2023 la consultora V\'ictor Gonz\'alez \& Asoc.~S.S.\ ejecut\'o un estudio
hidrogeol\'ogico para la Facultad de Ciencias y Tecnolog\'ias de la UC, Campus Asunci\'on
--- Sede Santa Librada~\cite{GonzalezHidro2023}, sobre un \'area de proyecto de
\SI{0.186}{\kilo\meter\squared}. El estudio incluy\'o tres \textbf{Sondeos El\'ectricos
Verticales} (SEV) con arreglo Schlumberger, medidos con un resistiv\'imetro GOLD~DDC-8
(hasta \SI{700}{\volt} y \SI{3}{\ampere}, $AB/2$ m\'aximo \SI{500}{\meter}) e
interpretados con el programa IPI2win de la Universidad Estatal de Mosc\'u.

\textbf{El SEV~01 se ejecut\'o en la cancha de f\'utbol de la UC} (UTM
$X=\num{435595}$, $Y=\num{7198920}$), que es el mismo predio en el que se realizaron las
campa\~nas de adquisici\'on de este proyecto. Es, por lo tanto, una medici\'on
independiente del mismo subsuelo, con un m\'etodo geof\'isico completamente distinto
---resistividad el\'ectrica en corriente continua en lugar de ondas superficiales--- y
anterior al proyecto, de modo que no puede haber sido influida por \'el.

\subsubsection{Modelos geoel\'ectricos obtenidos}

\begin{table}[h]
\centering\footnotesize
\caption{Modelos de capas invertidos de los tres SEV~\cite{GonzalezHidro2023}.
$\rho$: resistividad; $h$: espesor; $d$: profundidad a la base. El SEV~01 es el
co-localizado con las campa\~nas de este proyecto.}
\label{tab:sev}
\begin{tabular}{@{}llrrrl@{}}
\toprule
Sondeo & Capa & $\rho$ [\si{\ohm\meter}] & $h$ [\si{\meter}] & $d$ [\si{\meter}] & Interpretaci\'on \\
\midrule
\multirow{6}{*}{\shortstack[l]{\textbf{SEV~01}\\ Cancha de f\'utbol\\ $AB/2=\SI{125}{\meter}$\\ cota \SI{72}{\meter} s.n.m.}}
 & 1 & 1186  & 1{,}00 & 1{,}00  & Suelo residual seco, muy compactado \\
 & 2 & 15    & 1{,}06 & 2{,}06  & Nivel conductivo delgado (arcilloso) \\
 & 3 & 105   & 2{,}09 & 4{,}15  & Sedimentos areno-cuarzosos \\
 & 4 & 24{,}2 & 19{,}9 & \textbf{24{,}0} & Cementaci\'on arcillosa; \textbf{base del suelo} \\
 & 5 & 126   & 47{,}7 & \textbf{71{,}8} & Acu\'ifero Cuaternario/Pati\~no \\
 & 6 & 6{,}42 & --- & --- & Areniscas con matriz arcillosa \\
\midrule
\multirow{4}{*}{\shortstack[l]{SEV~02\\ Lindero norte\\ $AB/2=\SI{125}{\meter}$}}
 & 1 & 220   & 1{,}00 & 1{,}00  & Suelo residual seco \\
 & 2 & 36{,}6 & 12{,}0 & 13{,}0 & Sedimentos con mayor arcilla \\
 & 3 & 95{,}8 & 9{,}73 & 22{,}7 & Nivel cuarzoso \\
 & 4 & 30{,}8 & --- & --- & Basal arcilloso \\
\midrule
\multirow{4}{*}{\shortstack[l]{SEV~03 (param\'etrico)\\ Junto al pozo FT01\\ $AB/2=\SI{80}{\meter}$}}
 & 1 & 202   & 1{,}02 & 1{,}02 & Suelo residual seco \\
 & 2 & 117   & 7{,}34 & 8{,}36 & Areno-cuarzoso resistivo \\
 & 3 & 54{,}4 & 56{,}1 & 64{,}5 & Acu\'ifero \\
 & 4 & 3{,}64 & --- & --- & Basal arcilloso \\
\bottomrule
\end{tabular}
\end{table}

El SEV~03 es \textbf{param\'etrico}: se midi\'o junto a un pozo tubular profundo existente
(FT01, \SI{50}{\meter} de profundidad, en operaci\'on, conductividad el\'ectrica
\SI{130}{\micro\siemens\per\centi\meter}, pH~7{,}1), de modo que su modelo est\'a anclado
a litolog\'ia conocida. Eso da trazabilidad a la interpretaci\'on de los otros dos.

\subsubsection{Qu\'e aporta esto al proyecto}

\paragraph{a) Contrastes esperables en el perfil $V_S$.} El SEV~01 sit\'ua la
\textbf{base del suelo en \SI{24.0}{\meter}} y la base de la secuencia media en
\SI{71.8}{\meter}. Debe subrayarse una salvedad metodol\'ogica: \textbf{un contraste de
resistividad no es un contraste de rigidez}. La resistividad responde a porosidad,
salinidad del fluido y contenido de arcilla; $V_S$ responde a la rigidez al corte del
esqueleto. Coinciden a menudo, porque ambos siguen los cambios litol\'ogicos, pero no por
construcci\'on. Por eso las profundidades de la Tabla~\ref{tab:sev} se enuncian aqu\'i como
\emph{contrastes esperables} contra los cuales contrastar el perfil invertido, y no como
una predicci\'on que el MASW deba reproducir. Aun as\'i, el valor de la comparaci\'on es
alto: una interfaz a \SI{24}{\meter} cae de lleno dentro del objetivo de
\SI{50}{\meter} del proyecto, y es exactamente el tipo de rasgo que un perfil $V_S$
correctamente invertido deber\'ia resolver.

\paragraph{b) El nivel fre\'atico, con un n\'umero concreto.} El estudio ubica el nivel
del agua \textbf{entre \SI{10}{} y \SI{30}{\meter} de profundidad}, con la aclaraci\'on
expresa de que la campa\~na se realiz\'o en \'epoca de escasas precipitaciones, de modo
que se trata de una cota de estiaje y no del r\'egimen anual. Esto cierra un cabo suelto
del \S\ref{sec:no-unicidad}: all\'i se se\~nala que ignorar el nivel fre\'atico equivale a
asumir una raz\'on de Poisson incorrecta en las capas saturadas, donde el comportamiento
no drenado lleva $\nu\to0{,}5$. Ahora existe un rango concreto y con fuente para
parametrizar esa transici\'on en la inversi\'on ---y, m\'as importante, la advertencia de
que su profundidad es \textbf{estacional}, por lo que la fecha de cada campa\~na de este
proyecto es un dato relevante y no un detalle administrativo.

\paragraph{c) Marco litol\'ogico.} El sitio se asienta sobre el \textbf{Grupo Asunci\'on}
(Cret\'acico), relleno de una fosa tect\'onica tipo semi-graben: la Formaci\'on Pati\~no
en la base (fanglomerados, del orden de \SI{100}{\meter} de espesor) y la Formaci\'on
Yaguar\'on por encima (areniscas cuarzosas de grano fino a medio, friables, con matriz
arcillosa), que constituye el nivel productor del \textbf{Acu\'ifero Pati\~no}. Se trata
de sedimentos cl\'asticos poco cementados: es decir, exactamente el r\'egimen de suelo
blando a intermedio para el cual se asumieron velocidades de onda Rayleigh
$V_R\approx\SIrange{150}{250}{\meter\per\second}$ al derivar la banda objetivo en la
Ec.~\eqref{eq:geo-banda-objetivo}. La hip\'otesis que sostiene el requisito de banda del
proyecto queda as\'i respaldada por la geolog\'ia documentada del sitio y no s\'olo por un
valor de manual.

\paragraph{d) Un paralelo instructivo sobre la limitaci\'on de apertura.} El propio
estudio declara que la exploraci\'on qued\'o limitada a \SI{125}{\meter} de profundidad
``por limitaciones de espacio en \'area eminentemente urbana''. Es la misma clase de
restricci\'on ---no instrumental sino de sitio--- que limita la apertura del arreglo de
este proyecto, y vale citarla porque muestra que el problema no es exclusivo de un
prototipo de bajo costo: un equipo comercial operado por especialistas enfrent\'o la
misma barrera en el mismo predio.

\paragraph{e) Contraste con el criterio de Comina.} La geometr\'ia efectivamente
desplegada (apertura $L=\SI{40}{\meter}$, v\'ease m\'as abajo) admite longitudes de onda
resueltas del orden de la propia apertura. Frente al criterio
$\lambda_{\max}>\SI{60}{\meter}$ de Comina \emph{et al.}~\cite{Comina2011}
(\S\ref{sec:buenas-practicas}), \textbf{la apertura actual es insuficiente} para el
r\'egimen en que la no unicidad deja de dominar el error de $V_{S,30}$. Esto identifica
con precisi\'on el pr\'oximo cuello de botella del proyecto: \emph{no es la electr\'onica,
es la longitud del tendido}. Ampliar la apertura no requiere diseño nuevo, requiere m\'as
nodos ---que es justamente lo que la arquitectura inal\'ambrica fue pensada para
permitir.

\paragraph{f) Una validaci\'on invasiva al alcance de la mano.} El estudio
\textbf{recomienda perforar un pozo de \SI{60}{\meter} exactamente en el sector del
SEV~01}, es decir en la cancha, con perfilaje geof\'isico de pozo (resistividad, potencial
espont\'aneo y gamma natural). Si esa perforaci\'on llega a ejecutarse, el proyecto
tendr\'ia una \textbf{referencia invasiva directa en el mismo punto donde midi\'o},
reproduciendo a escala de un sitio el experimento de la segunda parte de
InterPACIFIC~\cite{Garofalo2016b}: contrastar un m\'etodo no invasivo contra perforaci\'on.
Es el mejor escenario de validaci\'on que este trabajo podr\'ia aspirar a tener y se
registra como l\'inea futura prioritaria (\S\ref{sec:futuros}).

\nota{Falta confirmar que las coordenadas registradas en los metadatos de las capturas de
``Canchita'' caen efectivamente dentro del entorno del SEV~01 (UTM 435595 / 7198920).
La identificaci\'on por nombre de sector (``cancha de f\'utbol de la UC'') es inequ\'ivoca,
pero conviene dejar asentada la distancia real entre el centro del tendido y el punto del
SEV, porque de ella depende cu\'an literal puede ser la comparaci\'on.}

\subsection{Campa\~nas de medici\'on en campo}

Se realizaron campa\~nas de adquisici\'on en varios sitios, con conjuntos de datos crudos
catalogados y versionados con Git~LFS y deduplicaci\'on por SHA-256. El cat\'alogo registra
del orden de \textbf{22\,000 archivos crudos} (aproximadamente \SI{28}{\gibi\byte})
distribuidos en tres sitios de medici\'on, m\'as los conjuntos procesados
correspondientes.

La geometr\'ia efectivamente usada fue un \textbf{arreglo virtual o apertura sint\'etica},
no un tendido f\'isico simult\'aneo de 21 ge\'ofonos. El martillo permaneci\'o en
$x=\SI{0}{\meter}$ y un ge\'ofono se desplaz\'o en incrementos de
$\Delta x=\SI{2}{\meter}$; cada registro contiene dos roles
(\texttt{n\_slaves=2}): martillo y ge\'ofono. El \emph{gather} se reconstruye alineando
capturas de impactos distintos mediante el canal de martillo.

Esta econom\'ia de hardware introduce una hip\'otesis adicional respecto del MASW
multicanal simult\'aneo: la fuente y el acoplamiento deben ser suficientemente repetibles,
y las condiciones del sitio deben permanecer estacionarias mientras se mueve el receptor.
Se registraron m\'ultiples repeticiones por posici\'on para apilar y evaluar esa
repetibilidad. En el conjunto procesado vigente, el grupo~1 cubre
$x=\SIrange{10}{50}{\meter}$ con 21 posiciones y apertura
$L=\SI{40}{\meter}$; el grupo~2 cubre $x=\SIrange{10}{36}{\meter}$ con 14 posiciones y
$L=\SI{26}{\meter}$.

\begin{figure}[ht]
\centering
\includegraphics[width=\textwidth]{captura_cruda_campo.png}
\caption{Registro crudo real usado para ilustrar el procedimiento de campo
($x=\SI{24}{\meter}$): canal de martillo, ge\'ofono con su nivel de continua y salida
FIR. La marca temporal proviene del impacto; una traza individual no demuestra
repetibilidad entre golpes ni coherencia equivalente a un arreglo simult\'aneo.}
\label{fig:captura-cruda-metodo}
\end{figure}

\subsection{Procesamiento de los registros de campo}

El tratamiento de los registros incorpor\'o varias correcciones sistem\'aticas que
surgieron de la propia pr\'actica:

\begin{itemize}[itemsep=0.15em]
    \item \textbf{Polaridad}: se estableci\'o una convenci\'on fija (ge\'ofono no invertido,
    martillo invertido) y se implement\'o correcci\'on de polaridad por punto de medici\'on,
    autom\'atica y manual, dado que un error de signo en un canal corrompe la imagen de
    dispersi\'on.
    \item \textbf{Alineaci\'on entre tandas (\emph{enfase})}: aplicada por carpeta de
    medici\'on, con una etapa de limpieza de desplazamientos por se\~nal.
    \item \textbf{Mezcla de frecuencias de muestreo}: los registros tomados a
    \SI{2929}{\hertz} y a \SI{1020}{\hertz} se combinan mediante remuestreo racional.
    \item \textbf{Deduplicaci\'on por captura} y \textbf{rechazo expl\'icito de carpetas}
    defectuosas, que las excluye de los promedios sin alterar su validez individual.
\end{itemize}

\begin{figure}[ht]
\centering
\includegraphics[width=\textwidth]{preprocesamiento_espectral.png}
\caption{Ejemplo de preprocesamiento de un barrido: ventana temporal, espectro y
estimaci\'on de la componente arm\'onica. Esta separaci\'on entre dato crudo y estimador es
la base de las curvas de respuesta experimentales.}
\label{fig:preprocesamiento-espectral}
\end{figure}

\nota{Para la versi\'on de 15 p\'aginas: conservar \S5.3 completo (es el experimento
publicable) y \S5.4.1 resumido a los defectos F1 y F9 en un p\'arrafo. \S5.2, \S5.5 y
\S5.6 pueden comprimirse fuertemente. La geometr\'ia debe conservarse: sin ella no se
puede interpretar la profundidad alcanzable ni la naturaleza virtual del arreglo.}
